\chapter{Principle of Virtual Forces}
\label{ch:pvf}

\section{Principle of virtual forces (PVF) and related principles or theorems}
The virtual force system should be ``statically admissible'' (the virtual displacements should be ``kinematically admissible''); this means that the external virtual forces and the internal forces are in equilibrium.

For an admissible virtual force system, the principle of virtual forces is stated as
\begin{equation}
  \dW^* = \dU_i^* \quad\Longleftrightarrow\quad \text{Compatibility Equations}
\end{equation}
where $\dW^* =$ the complementary virtual work done by the external virtual forces through the ``real displacements'' (or the external complementary virtual work), and $\dU_i^* =$ the complementary virtual internal energy (or the internal complementary virtual work). They are defined, respectively, by
\begin{align}
  \dW^* &= \int_S\delta t_iu_i\,dS + \int_V\delta f_iu_i\,dV\\
  \dU_i^* &= \int_V\delta\sigma_{ij}\eps_{ij}\,dV
\end{align}
The statically admissible conditions consist of the following two equilibrium conditions.
\begin{gather*}
  \delta\sigma_{ij,j} + \delta f_i = 0\ \text{ in } V\\
  \delta\sigma_{ij}n_j = \delta t_i\ \text{ on } S\ (\text{not } S_t)
\end{gather*}
For an (linearly or nonlinearly) elastic body, there exists a function of stresses (internal forces), called complementary strain energy $U^*$, such that $\dU^* = \dU_i^*$. If the external forces are conservative, then there exists a complementary load potential $V^*$ such that
\begin{equation*}
  \delta V^* = -\dW^*
\end{equation*}
Therefore, for a conservative system because the material is elastic and the external forces are conservative, there exists the complementary total potential energy defined by
\begin{equation*}
  \Pi^* = U^* + V^*
\end{equation*}
The PVF then becomes
\begin{equation}
  \delta\Pi^* = 0 \quad\Longleftrightarrow\quad \text{Compatibility Equations}
\end{equation}
which is often referred to as the \emph{principle of stationary (minimum) total complementary potential energy}.

If the external forces acting on a conservative system are point (concentrated) loads, then
\begin{equation*}
  \Pi^* = U^* + V^* = U^*(\Pvec) - \Pvec^\tr\bs{\Delta} = U^*(\Pvec) - \sum_{i=1}^nP_i\Delta_i
\end{equation*}
Therefore,
\begin{equation*}
  \delta\Pi^* = 0 \;\Rightarrow\; \sum_{i=1}^n\Bigl(\pd{U^*}{P_i} - \Delta_i\Bigr)\delta P_i = 0
\end{equation*}
As $\delta P_i$'s are independent, we can conclude that
\begin{equation}
  \pd{U^*}{P_i} = \Delta_i,\quad i = 1, n
\end{equation}
This is often referred to as \emph{Engesser's first theorem}. Note that for \emph{linearly} elastic bodies, $U^*$ is a \emph{quadratic} function of $\Pvec$ and its value is the same as $U$. Therefore, in this case we often use the notation $U(\Pvec)$, i.e.\ the strain energy, instead of $U^*(\Pvec)$. The above relation then becomes
\begin{equation}
  \pd{U}{P_i} = \Delta_i,\quad i = 1, n
\end{equation}
This is often referred to as \emph{Castigliano's second theorem}.

For a statically indeterminate conservative structural system, the total complementary potential energy can be written as
\begin{equation*}
  \Pi^* = U^* + V^* = U^*(\Pvec, \mat{R}) - \Pvec^\tr\bs{\Delta} - \mat{R}^\tr\mat{r} = U^*(\Pvec, \mat{R}) - \sum_{i=1}^nP_i\Delta_i - \sum_{i=1}^rR_ir_i
\end{equation*}
The principle of stationary (minimum) total complementary potential energy states that
\begin{equation}
  \pd{U^*}{P_i} = \Delta_i,\quad i = 1, n
\end{equation}
and
\begin{equation}
  \pd{U^*}{R_i} = r_i = 0,\quad i = 1, r
\end{equation}
The condition $r_i = 0$ is very obvious if the associated $R_i$ is a reaction because no movement is allowed at the corresponding support. If $R_i$ is an internal force, the condition means consistency of deformation (no gap, continuous slope, etc.). \emph{The second set of equations is often referred to as the theorem of least work or Engesser's second theorem.}

In matrix form, the above two sets of equations can be written as
\begin{equation}
  \Bmat{\dfrac{\partial U^*}{\partial\Pvec}\\[8pt] \dfrac{\partial U^*}{\partial\mat{R}}} = \Bmat{\bs{\Delta}\\ \mat{0}}
\end{equation}
Having obtained the above equations, we can solve for $\mat{R}$ from the second set of equations. Finally, substitute $\mat{R}$ into the first set of equations to obtain $\bs{\Delta}$.

\section{Complementary strain energy and the unit-load method}
The complementary strain energy for a \emph{linearly} elastic structural member:
\begin{enumerate}[label=(\alph*)]
  \item due to axial deformation,
    \begin{equation*}
      U^* = \int\frac{N^2}{2EA}\,dx = \frac{N^2L}{2EA}\ \text{ if } \frac{N}{EA} = \text{constant, where } N = \text{axial force.}
    \end{equation*}
  \item due to flexural deformation,
    \begin{equation*}
      U^* = \int\frac{M^2}{2EI}\,dx,\ \text{ where } M = \text{bending moment}
    \end{equation*}
  \item due to torsional deformation,
    \begin{equation*}
      U^* = \int\frac{T^2}{2GJ}\,dx = \frac{T^2L}{2GJ}\ \text{ if } \frac{T}{GJ} = \text{constant, where } T = \text{torsion.}
    \end{equation*}
\end{enumerate}
The unit-load method is an application of the principle of virtual forces. Consider a structural system consisting of \emph{linearly elastic} members, for which we like to find the displacement $\Delta_i$. Define a virtual force system that consists of only one virtual force $\delta P_i$ acting on the structural system and along the direction of $\Delta_i$. According to the PVF,
\begin{equation}
  \delta P_i\Delta_i = \int\frac{Nn}{EA}\,dx + \int\frac{Mm}{EI}\,dx + \int\frac{Tt}{GJ}\,dx
\end{equation}
where $N$, $M$, and $T =$ axial force, bending moment, and torsion due to the real external loads, and $n$, $m$, and $t =$ axial force, bending moment, and torsion due to the virtual force $\delta P_i$. As the virtual internal forces $n$, $m$, and $t$ depend on $\delta P_i$ linearly, without loss of generality we could let $\delta P_i = 1$; this results in
\begin{equation}
  \Delta_i = \int\frac{Nn}{EA}\,dx + \int\frac{Mm}{EI}\,dx + \int\frac{Tt}{GJ}\,dx
\end{equation}
where $n$, $m$, and $t =$ axial force, bending moment, and torsion due to $\delta P_i = 1$. This is the so-called \emph{unit load method}.
\begin{fig}
  % Real system (N, M, T) and virtual system (n, m, t); support conditions can be different
\begin{tikzpicture}[scale=0.9,font=\small]
  \begin{scope}
    \draw[thick] plot[smooth cycle,tension=0.9] coordinates {(0,0) (1.5,-0.9) (3.2,-0.3) (3.4,1.2) (2.0,2.0) (0.4,1.4)};
    \draw[load] (-1.1,1.2) -- (0.4,0.6) node[pos=0,above] {$P_1$};
    \draw[load] (1.6,0.4) -- (1.6,2.4) node[above] {$P_2$};
    \draw[load] (2.2,1.0) -- (3.9,1.0) node[right] {$P_3$};
    \pic[rotate=-30] at (3.3,0.2) {pin}; \pic[rotate=20] at (1.3,-0.9) {roller};
    \node[below] at (1.6,-1.6) {\itshape real $(N, M, T)$};
  \end{scope}
  \begin{scope}[xshift=7cm]
    \draw[thick] plot[smooth cycle,tension=0.9] coordinates {(0,0) (1.5,-0.9) (3.2,-0.3) (3.4,1.2) (2.0,2.0) (0.4,1.4)};
    \draw[load] (-0.4,2.2) -- (0.8,1.0) node[pos=0,above] {$\delta P$};
    \pic[rotate=-30] at (3.3,1.2) {pin}; \pic at (1.5,-0.9) {roller};
    \node[below] at (1.6,-1.6) {\itshape virtual $(n, m, t)$};
  \end{scope}
\end{tikzpicture}

\end{fig}
\noindent$***$ Support conditions can be different!

\section{Nonlinear system}
How can we use the unit load method to deal with nonlinear systems? Assume that
\begin{equation*}
  U^* = \int\frac{aM^b}{EI}\,dx
\end{equation*}
where $a$ and $b$ are constants. Then, $\dU^* = \int\dfrac{abM^{b-1}}{EI}\delta M\,dx$. Therefore the unit-load method takes the following form:
\begin{equation}
  \Delta_i = \int\frac{abM^{b-1}}{EI}m\,dx.
\end{equation}

\begin{example}[Application of P.V.F.]
\leavevmode
\begin{fig}
  % Real system (cantilever with tip load) and virtual system 1 (simply supported, load at C)
\begin{tikzpicture}[scale=0.9,font=\small]
  \begin{scope}
    \pic at (0,0) {wallL}; \draw[member] (0,0) -- (6,0);
    \node[above right] at (0,0.05) {$A$}; \node[above] at (3,0.05) {$C$}; \node[above left] at (6,0.05) {$B$};
    \draw[load] (6,0.1) -- (6,1.3) node[above] {$P$};
    \draw[dim] (0,-0.6) -- (3,-0.6) node[midway,fill=white,inner sep=1pt] {$\frac L2$}; \draw[dim] (3,-0.6) -- (6,-0.6) node[midway,fill=white,inner sep=1pt] {$\frac L2$};
    \node[below] at (3,-1.0) {\itshape real system};
  \end{scope}
  \begin{scope}[xshift=8.5cm]
    \draw[member] (0,0) -- (6,0); \pic at (0,0) {roller}; \pic at (6,0) {pin}; \node[node dot] at (0,0) {}; \node[node dot] at (6,0) {};
    \node[above] at (0.3,0.05) {$A$}; \node[above right] at (3,0.05) {$C$}; \node[above] at (5.7,0.05) {$B$};
    \draw[load] (3,0.1) -- (3,1.3) node[above] {$\delta P$};
    \draw[dim] (0,-0.8) -- (3,-0.8) node[midway,fill=white,inner sep=1pt] {$\frac L2$}; \draw[dim] (3,-0.8) -- (6,-0.8) node[midway,fill=white,inner sep=1pt] {$\frac L2$};
    \node[below] at (3,-1.2) {\itshape virtual system 1};
  \end{scope}
\end{tikzpicture}

\end{fig}
For real system: $M(x) = P(L-x)$

\emph{Case 1: virtual system 1}

\emph{The boundary conditions for the real system and the virtual system are totally different, but the P.V.F.\ still works as we shall show below.}

The virtual internal moment distribution is
\begin{align*}
  m(x) &= -\frac{\delta P}{2}x,\quad 0 < x < \frac L2\\
       &= -\frac{\delta P}{2}(L-x),\quad \frac L2 < x < L
\end{align*}
The complementary virtual work and the complementary virtual strain energy are given respectively by
\begin{gather*}
  \dW^* = -\frac{\delta P}{2}\times0 + \delta P\times v_c - \frac{\delta P}{2}\times v_B\\
  \begin{aligned}
    \dU^* &= \int_0^L\frac{M(x)m(x)}{EI}\,dx\\
          &= \frac{1}{EI}\int_0^{L/2}\Bigl[P(L-x)\Bigl(-\frac{\delta P}{2}x\Bigr)\Bigr]dx + \frac{1}{EI}\int_{L/2}^LP(L-x)\Bigl(-\frac{\delta P}{2}\Bigr)(L-x)\,dx\\
          &= -\frac{\delta P}{16}\frac{PL^3}{EI}
  \end{aligned}
\end{gather*}
Then, by the P.V.F.\ $\dW^* = \dU^*$, the following ``compatibility equation'' can be obtained.
\begin{equation*}
  v_C - \frac{v_B}{2} = -\frac{1}{16}\frac{PL^3}{EI}
\end{equation*}
Although this relation may not be useful, it is valid since the exact values are $v_B = \dfrac{PL^3}{3EI}$ and $v_C = \dfrac{5PL^3}{48EI}$.

\emph{Case 2: virtual system 2}
\begin{fig}
  % Virtual system 2: fixed at B, load at A
\begin{tikzpicture}[scale=0.9,font=\small]
  \draw[member] (0,0) -- (6,0); \pic at (6,0) {wallR};
  \node[above right] at (0,0.05) {$A$}; \node[above] at (3,0.05) {$C$}; \node[above left] at (6,0.05) {$B$};
  \draw[load] (0,0.1) -- (0,1.3) node[above] {$\delta P$};
  \draw[dim] (0,-0.6) -- (3,-0.6) node[midway,fill=white,inner sep=1pt] {$\frac L2$}; \draw[dim] (3,-0.6) -- (6,-0.6) node[midway,fill=white,inner sep=1pt] {$\frac L2$};
  \node[below] at (3,-1.0) {\itshape virtual system 2};
\end{tikzpicture}

\end{fig}
\begin{gather*}
  m(x) = \delta P\cdot x\\
  \dW^* = \delta P\cdot L\times\theta_B - \delta P\times v_B\\
  \dU^* = \int_0^L\frac{M(x)m(x)}{EI}\,dx = \frac{\delta P}{6}\frac{PL^3}{EI}
\end{gather*}
Then, by P.V.F.\ $\dW^* = \dU^*$, the following compatibility equation can be obtained.
\begin{equation*}
  -v_B + \theta_BL = \frac16\frac{PL^3}{EI}
\end{equation*}
which is true (but may not be so much useful) since $v_B = \dfrac{PL^3}{3EI}$ and $\theta_B = \dfrac{PL^2}{2EI}$.

\emph{Case 3: virtual system 3: solve the vertical displacement of $B$}
\begin{fig}
  % Virtual system 3: fixed at A, load at B
\begin{tikzpicture}[scale=0.9,font=\small]
  \pic at (0,0) {wallL}; \draw[member] (0,0) -- (6,0);
  \node[above right] at (0,0.05) {$A$}; \node[above] at (3,0.05) {$C$}; \node[above left] at (6,0.05) {$B$};
  \draw[load] (6,0.1) -- (6,1.3) node[above] {$\delta P$};
  \draw[dim] (0,-0.6) -- (3,-0.6) node[midway,fill=white,inner sep=1pt] {$\frac L2$}; \draw[dim] (3,-0.6) -- (6,-0.6) node[midway,fill=white,inner sep=1pt] {$\frac L2$};
  \node[below] at (3,-1.0) {\itshape virtual system 3};
\end{tikzpicture}

\end{fig}
\begin{gather*}
  m(x) = \delta P(L-x)\\
  \dW^* = \delta P\times v_B\\
  \dU^* = \int_0^L\frac{M(x)m(x)}{EI}\,dx = \frac{\delta P}{3}\frac{PL^3}{EI}
\end{gather*}
Then, by P.V.F.\ $\dW^* = \dU^*$, the vertical displacement of $B$ can be obtained as
\begin{equation*}
  v_B = \frac13\frac{PL^3}{EI}
\end{equation*}

\emph{Case 4: virtual system 4: solve the vertical displacement of $C$}
\begin{fig}
  % Virtual system 4: fixed at A, load at C
\begin{tikzpicture}[scale=0.9,font=\small]
  \pic at (0,0) {wallL}; \draw[member] (0,0) -- (6,0);
  \node[above right] at (0,0.05) {$A$}; \node[above right] at (3,0.05) {$C$}; \node[above left] at (6,0.05) {$B$};
  \draw[load] (3,0.1) -- (3,1.3) node[above] {$\delta P$};
  \draw[dim] (0,-0.6) -- (3,-0.6) node[midway,fill=white,inner sep=1pt] {$\frac L2$}; \draw[dim] (3,-0.6) -- (6,-0.6) node[midway,fill=white,inner sep=1pt] {$\frac L2$};
  \node[below] at (3,-1.0) {\itshape virtual system 4};
\end{tikzpicture}

\end{fig}
The virtual force system shall be chosen as that shown.

Note that the following formula is very useful when carrying out the integration of $\int_0^LM(x)m(x)\,dx$. Some other specific cases can be seen in the last page of this chapter (Table~\ref{tab:integral}), which can be considered special cases of this formula.
\begin{fig}
  % Simpson-type formula for the integral of M(x) m(x)
\begin{tikzpicture}[scale=0.9,font=\small]
  \begin{scope}
    \draw (0,0) -- (5,0);
    \draw[thick] (0,0) -- (0,0.9) .. controls (1.5,1.6) and (3.5,1.9) .. (5,2.0) -- (5,0);
    \draw (2.5,0) -- (2.5,1.63);
    \node[right] at (0.05,0.6) {$a_1$}; \node[right] at (2.55,1.2) {$a_2$}; \node[right] at (5.05,1.6) {$a_3$};
    \node[left] at (-0.1,0.5) {$M(x)$};
    \node[below] at (0,-0.05) {0}; \node[below] at (2.5,-0.05) {$\frac L2$}; \node[below] at (5,-0.05) {$L$};
  \end{scope}
  \begin{scope}[xshift=7.5cm]
    \draw (0,0) -- (5,0);
    \draw[thick] (0,0) -- (0,0.5) -- (5,2.2) -- (5,0);
    \draw (2.5,0) -- (2.5,1.35);
    \node[right] at (0.05,0.3) {$b_1$}; \node[right] at (2.55,1.0) {$b_2$}; \node[right] at (5.05,1.8) {$b_3$};
    \node[left] at (-0.1,0.5) {$m(x)$};
    \node[below] at (0,-0.05) {0}; \node[below] at (2.5,-0.05) {$\frac L2$}; \node[below] at (5,-0.05) {$L$};
  \end{scope}
\end{tikzpicture}

\end{fig}
\begin{equation}
  \int_0^LM(x)m(x)\,dx = \frac L6\bigl(a_1b_1 + 4a_2b_2 + a_3b_3\bigr)
\end{equation}
\end{example}

\begin{example}[Application of the unit load method]
\leavevmode

\begin{minipage}{0.55\linewidth}\centering% Propped cantilever with an overhang: 30 kN at the tip
\begin{tikzpicture}[scale=0.9,font=\small]
  \pic at (0,0) {wallL}; \draw[member] (0,0) -- (9,0); \pic at (6,0) {roller};
  \node[above right] at (0,0.05) {$A$}; \node[above] at (6,0.05) {$C$}; \node[above left] at (9,0.05) {$B$};
  \draw[load] (9,0.1) -- (9,1.4) node[above] {30 kN};
  \draw[dim] (0,-1.0) -- (6,-1.0) node[midway,fill=white,inner sep=1pt] {6 m}; \draw[dim] (6,-1.0) -- (9,-1.0) node[midway,fill=white,inner sep=1pt] {3 m};
  \foreach \x in {0,6,9}{\draw (\x,-0.8) -- (\x,-1.2);}
\end{tikzpicture}
\end{minipage}\hfill
\begin{minipage}{0.42\linewidth}
  $EI =$ constant
  \begin{align*}
    M(x) &= -45 + ax,\quad 0\leq x\leq 6\\
         &= c - bx,\quad 6\leq x\leq 9
  \end{align*}
\end{minipage}
\begin{enumerate}[label=(\alph*)]
  \item Determine the values of the coefficients $a$, $b$ and $c$.
  \item Determine $\theta_B$.
\end{enumerate}
\begin{solution}
\begin{enumerate}[label=(\alph*)]
  \item The value of $b = 30$ is obvious since the slope in the moment diagram is equal to the shear force. Knowing that $b = 30$ and $M(9) = 0$, we obtain $c = 270$. Finally, the moment at $x = 6$ shall be continuous, i.e.
    \begin{equation*}
      M(6) = c - bx = 270 - 30\times6 = 90 = -45 + 6a
    \end{equation*}
    Therefore, $a = 22.5$.
  \item Since the virtual system can be of any kind of boundary conditions, the simplest one (a cantilever) is
    \begin{fig}
      % Virtual system: cantilever with a unit moment at B
\begin{tikzpicture}[scale=0.9,font=\small]
  \pic at (0,0) {wallL}; \draw[member] (0,0) -- (9,0);
  \node[above right] at (0,0.05) {$A$}; \node[above] at (6,0.05) {$C$}; \node[above left] at (9,0.05) {$B$};
  \draw[->,thick] (9.5,0.5) arc (60:-120:0.5); \node[right] at (9.7,0) {$m = 1$};
  \draw[dim] (0,-1.0) -- (6,-1.0) node[midway,fill=white,inner sep=1pt] {6 m}; \draw[dim] (6,-1.0) -- (9,-1.0) node[midway,fill=white,inner sep=1pt] {3 m};
  \foreach \x in {0,6,9}{\draw (\x,-0.8) -- (\x,-1.2);}
\end{tikzpicture}

    \end{fig}
    The virtual moment distribution is $m(x) = 1$. Then, $\dW^* = \dU^*$ becomes
    \begin{equation*}
      1\times\theta_B = \frac{1}{EI}\int_0^6(-45 + 22.5x)\,dx + \frac{1}{EI}\int_6^9(270 - 30x)\,dx = \frac{270}{EI}
    \end{equation*}
\end{enumerate}
\end{solution}
\end{example}

\begin{example}[Application of the unit load method]
\leavevmode
\begin{fig}
  % Propped cantilever with a central load: real system (redundant R_B) and virtual system
\begin{tikzpicture}[scale=0.9,font=\small]
  \begin{scope}
    \pic at (0,0) {wallL}; \draw[member] (0,0) -- (6,0); \pic at (6,0) {roller}; \node[node dot] at (6,0) {};
    \node[above right] at (0,0.05) {$A$}; \node[above left] at (3,0.05) {$C$}; \node[above left] at (6,0.05) {$B$};
    \draw[load] (3,1.3) -- (3,0.1) node[pos=0,above] {$P$};
    \draw[load] (6,-1.6) -- (6,-0.7) node[pos=0,below] {$R_B$};
    \draw[dim] (0,-1.1) -- (3,-1.1) node[midway,fill=white,inner sep=1pt] {$\frac L2$}; \draw[dim] (3,-1.1) -- (6,-1.1) node[midway,fill=white,inner sep=1pt] {$\frac L2$};
    \node[below] at (2.5,-1.6) {\itshape real system};
  \end{scope}
  \begin{scope}[xshift=8.5cm]
    \pic at (0,0) {wallL}; \draw[member] (0,0) -- (6,0);
    \node[above right] at (0,0.05) {$A$}; \node[above] at (3,0.05) {$C$}; \node[above left] at (6,0.05) {$B$};
    \draw[load] (6,0.1) -- (6,1.3) node[above] {1};
    \draw[dim] (0,-1.1) -- (3,-1.1) node[midway,fill=white,inner sep=1pt] {$\frac L2$}; \draw[dim] (3,-1.1) -- (6,-1.1) node[midway,fill=white,inner sep=1pt] {$\frac L2$};
    \node[below] at (3,-1.6) {\itshape virtual system};
  \end{scope}
\end{tikzpicture}

\end{fig}
Take the reaction at the roller as the redundant force. Then the moment distribution for the real system is
\begin{align*}
  M(x) &= R_B(L-x),\quad \frac L2 < x < L\\
       &= R_B(L-x) - P\Bigl(\frac L2 - x\Bigr),\quad 0 < x < \frac L2
\end{align*}
From the virtual system, we have $m(x) = 1\times(L-x)$. According to the unit load method
\begin{equation*}
  v_B = \int_0^L\frac{Mm}{EI}\,dx = 0 \;\Rightarrow\; 0 = \frac{5PL^3}{48EI} - \frac{R_BL^3}{3EI} \;\Rightarrow\; R_B = \frac{5}{16}P
\end{equation*}
\begin{fig}
  % Reactions of the propped cantilever
\begin{tikzpicture}[scale=0.9,font=\small]
  \pic at (0,0) {wallL}; \draw[member] (0,0) -- (6,0); \pic at (6,0) {pin}; \node[node dot] at (6,0) {};
  \draw[load] (3,1.3) -- (3,0.1) node[pos=0,above] {$P$};
  \draw[load] (0,-1.6) -- (0,-0.5) node[pos=0,below] {$\frac{11}{16}P$};
  \draw[load] (6,-1.6) -- (6,-0.7) node[pos=0,below] {$\frac{5}{16}P$};
  \draw[->,thick] (-0.9,0.6) arc (110:-40:0.5); \node[left] at (-0.9,0.5) {$\frac{PL}{8}\times\frac32$};
\end{tikzpicture}

\end{fig}
\end{example}

\begin{example}[Application of Engesser's 1st theorem]
\leavevmode
\begin{fig}
  % Truss for Engesser's first theorem and the nonlinear member response
\begin{tikzpicture}[scale=0.9,font=\small]
  \begin{scope}
    \coordinate (a) at (0,0); \coordinate (c) at (3.2,0); \coordinate (b) at (1.6,-2.77); \coordinate (d) at (4.8,-2.77);
    \draw[member] (a) -- (c) -- (b) -- (a); \draw[member] (c) -- (d);
    \pic at (b) {pin}; \pic at (d) {pin};
    \node[above] at (a) {$a$}; \node[above] at (c) {$c$}; \node[below left] at (b) {$b$}; \node[below right] at (d) {$d$};
    \draw[load] (a) -- ++(-1.4,0) node[left] {$P$};
    \node[above] at (1.6,0.05) {$(P)$}; \node[left] at (0.6,-1.4) {$(0)$}; \node[right] at (2.5,-1.4) {$(-P)$}; \node[right] at (4.1,-1.4) {$(P)$};
    \draw[->] ($(a)+(1.0,0)$) arc (0:-60:1.0); \node[right] at ($(a)+(0.9,-0.6)$) {$60^\circ$};
  \end{scope}
  \begin{scope}[xshift=8cm,yshift=-2.8cm]
    \draw[->] (0,0) -- (5.2,0) node[right] {$\delta$}; \draw[->] (0,0) -- (0,4.2) node[above] {$f$};
    \draw[thick,domain=0:4.6,samples=40] plot (\x,{2.3*sqrt(\x/1.0)*0.85});
    \pgfmathsetmacro{\F}{2.3*sqrt(1.4)*0.85}
    \fill[pattern={Lines[angle=45,distance=3pt]}] (0,0) -- plot[domain=0:1.4,samples=30] (\x,{2.3*sqrt(\x)*0.85}) -- (0,\F) -- cycle;
    \draw[dashed] (0,\F) -- (1.4,\F); \node[left] at (0,\F) {$F$};
    \draw[dashed] (0,{\F*0.62}) -- (0.55,{\F*0.62}); \draw[dashed] (0,{\F*0.62+0.25}) -- (0.62,{\F*0.62+0.25});
    \draw[dim] (-0.25,{\F*0.62}) -- (-0.25,{\F*0.62+0.25}); \node[left] at (-0.3,{\F*0.62+0.12}) {$df$};
    \draw[dim] (0,-0.3) -- (0.55,-0.3); \node[below] at (0.3,-0.35) {$\delta$};
    \node[right] at (1.8,1.6) {$\delta = K_1f + K_2f^2$};
  \end{scope}
\end{tikzpicture}

\end{fig}
The complementary strain energy for a member:
\begin{equation*}
  U^*(F) = \int_0^F\delta(f)\,df = \int_0^F(K_1f + K_2f^2)\,df = \tfrac12K_1F^2 + \tfrac13K_2F^3\quad\text{(the shaded area)}
\end{equation*}
The structure is statically determinate and the member forces are already shown in the left figure. The complementary strain energy of the system is
\begin{equation*}
  U^* = \sum U^*_j = \tfrac12K_1\bigl(P^2 + (-P)^2 + P^2\bigr) + \tfrac13K_2\bigl(P^3 + (-P)^3 + P^3\bigr) = \tfrac32K_1P^2 + \tfrac13K_2P^3
\end{equation*}
By Engesser's 1st theorem, the horizontal displacement ($\leftarrow$) at the loaded point is
\begin{equation*}
  \Delta_a = \pd{U^*}{P} = 3K_1P + K_2P^2\quad(\leftarrow)
\end{equation*}
\end{example}

\begin{example}[3D structure]
\leavevmode
\begin{fig}
  % 3D L-shaped structure: cantilever AB with a perpendicular arm BC loaded by P at C
\begin{tikzpicture}[scale=1.0,font=\small,x={(1cm,0cm)},y={(0cm,1cm)},z={(-0.55cm,-0.85cm)}]
  \begin{scope}
    \def\h{0.5} \def\w{0.5}
    % member AB along +x from A(0,0,0) to B(3,0,0)
    \draw (0,0,0) -- (3,0,0) -- (3,\h,0) -- (0,\h,0) -- cycle;
    \draw (3,0,0) -- (3,0,\w) -- (3,\h,\w) -- (3,\h,0); \draw (0,\h,0) -- (0,\h,\w) -- (3,\h,\w);
    % member BC along +z from B to C(3,0,4.5)
    \draw (3,0,\w) -- (3,0,4.5) -- (3,\h,4.5) -- (3,\h,\w); \draw (3,\h,4.5) -- (2.5,\h,4.5) -- (2.5,0,4.5) -- (3,0,4.5);
    \draw (2.5,\h,4.5) -- (2.5,\h,\w);
    % wall at A
    \fill[hatch] (-0.25,-0.3,0) -- (0,-0.3,0) -- (0,\h+0.3,0) -- (-0.25,\h+0.3,0) -- cycle;
    \node[above left] at (0,\h,0) {$A$}; \node[above] at (3,\h,0) {$B$}; \node[below] at (2.75,0,4.6) {$C$};
    \draw[->] (2.2,\h+0.15,0) -- (1.0,\h+0.15,0) node[above] {$x'$};
    \draw[->] (2.75,\h+0.1,3.0) -- (2.75,\h+0.1,1.8) node[above right] {$x'$};
    \draw[load] (2.75,\h+1.5,4.2) -- (2.75,\h+0.15,4.2) node[pos=0,left] {$P$};
    \draw[dim] (0,\h+0.9,0) -- (3,\h+0.9,0) node[midway,above] {$L$};
    \draw[dim] (3.6,0,\w) -- (3.6,0,4.5) node[midway,right] {$2L$};
  \end{scope}
  % section
  \begin{scope}[shift={(7.5,0.5,0)}]
    \draw (-0.5,-0.5,0) rectangle (0.5,0.5,0); \draw (-0.35,-0.35,0) rectangle (0.35,0.35,0);
    \draw[->,thick] (0,0,0) -- (0,2.0,0) node[above] {$I_2\,(z')$};
    \draw[->,thick] (0,0,0) -- (2.0,0,0) node[right] {$I_1\,(y')$};
    \draw[->,thick] (0,0,0) -- (0,0,2.2) node[below] {$J\,(x')$};
    \node[right] at (0.4,-1.0,0) {section};
  \end{scope}
\end{tikzpicture}

\end{fig}
Determine the vertical displacement at node $C$.
\begin{solution}
\leavevmode
\begin{fig}
  % Free bodies of members AB and BC
\begin{tikzpicture}[scale=0.9,font=\small]
  \begin{scope}
    \node[left] at (-0.8,0) {Member $AB$:};
    \pic at (0,0) {wallL}; \draw (0,-0.25) rectangle (5,0.25);
    \node[above right] at (0,0.25) {$A$}; \node[above left] at (5,0.25) {$B$};
    \draw[->] (4.3,0) -- (2.8,0) node[below] {$x'$};
    \draw[load] (5,1.4) -- (5,0.3) node[pos=0,above] {$P$};
    \draw[->,thick,double] (5.2,0) -- (6.4,0) node[right] {$2PL$};
    \draw[->] (0,-0.3) -- (0,-1.5) node[below] {$z'$};
  \end{scope}
  \begin{scope}[yshift=-4.8cm]
    \node[left] at (-0.8,0) {Member $BC$:};
    \draw (0,-0.25) rectangle (7,0.25);
    \node[above right] at (0,0.25) {$C$}; \node[above left] at (7,0.25) {$B$};
    \draw[->] (0.2,0) -- (1.7,0) node[below] {$x'$};
    \draw[load] (0,1.4) -- (0,0.3) node[pos=0,above] {$P$};
    \draw[->] (0,-0.3) -- (0,-1.5) node[below] {$z'$};
  \end{scope}
\end{tikzpicture}

\end{fig}
Member $AB$: $M(x') = Px'$ (w.r.t.\ the $y'$ axis), $T(x') = 2PL$.\\
Member $BC$: $M(x') = Px'$ (w.r.t.\ the $y'$ axis), $T(x') = 0$.

According to the unit load method
\begin{equation*}
  1\cdot v_C = \int_0^L\frac{(Px)x}{EI_1}\,dx + \int_0^L\frac{(2PL)2L}{GJ}\,dx + \int_0^{2L}\frac{(Px)x}{EI_1}\,dx = 3\frac{PL^3}{EI_1} + 4\frac{PL^3}{GJ}
\end{equation*}
\end{solution}
\end{example}

\begin{example}[Deflection of truss due to external loads]
Calculate the horizontal displacement of node $C$, $D_1$, and the \emph{relative displacement} of nodes $B$ and $E$, $D_2$. Note that the cross section for members 6 and 7 $= 1.25A$ and that for other members is $A$. (Ghali and Neville 1997, \emph{Structural Analysis})
\begin{fig}
  % Truss example (Ghali & Neville): geometry, real forces, and virtual forces
\begin{tikzpicture}[scale=0.62,font=\scriptsize]
  \newcommand{\trussgeo}{\coordinate (A) at (0,0); \coordinate (B) at (4,0); \coordinate (C) at (8,0); \coordinate (D) at (4,3); \coordinate (E) at (0,3);
    \draw[thinmember] (A) -- (C); \draw[thinmember] (E) -- (D) -- (C); \draw[thinmember] (A) -- (E); \draw[thinmember] (A) -- (D); \draw[thinmember] (B) -- (D);
    \foreach \p in {A,B,C,D,E}{\fill (\p) circle (0.06);}}
  % (a)
  \begin{scope}
    \trussgeo \pic at (A) {pin}; \pic at (C) {roller};
    \node[left] at (A) {$A$}; \node[below] at (B) {$B$}; \node[below right] at (C) {$C$}; \node[above right] at (D) {$D$}; \node[above left] at (E) {$E$};
    \newcommand{\mem}[2]{\node[circle,draw,fill=white,inner sep=0.8pt,font=\tiny] at #1 {#2};}
    \mem{(2,0)}{1} \mem{(6,0)}{2} \mem{(0,1.5)}{3} \mem{(4,1.5)}{4} \mem{(2,3)}{5} \mem{(6,1.5)}{6} \mem{(1.5,1.4)}{7}
    \draw[load] (-1.3,3) -- (-0.1,3) node[pos=0,above] {$P$}; \draw[load] (4,4.3) -- (4,3.1) node[pos=0,right] {$P$};
    \draw[load] (8.1,0) -- (9.1,0) node[right] {$D_1$}; \node[below] at (8.6,0) {1};
    \draw[load] (E) -- ++(-40:0.9) node[below right] {$\frac{D_2}{2}$}; \draw[load] (B) -- ++(140:0.9) node[above left] {$\frac{D_2}{2}$};
    \draw[dim] (0,-1.0) -- (4,-1.0) node[midway,fill=white,inner sep=1pt] {$l$}; \draw[dim] (4,-1.0) -- (8,-1.0) node[midway,fill=white,inner sep=1pt] {$l$};
    \node[below] at (4,-1.5) {(a) 2D truss};
  \end{scope}
  % (b)
  \begin{scope}[xshift=12cm]
    \trussgeo \pic at (A) {pin}; \pic at (C) {roller};
    \draw[load] (-1.3,3) -- (-0.1,3) node[pos=0,above] {$P$}; \draw[load] (4,4.3) -- (4,3.1) node[pos=0,right] {$P$};
    \draw[load] (-1.3,0) -- (-0.1,0) node[pos=0,above] {$P$};
    \node[above] at (2,3) {$-P$}; \node[left] at (0,1.5) {$0$}; \node[right] at (4,1.5) {$0$}; \node[below] at (2,0) {$1.167P$}; \node[below] at (6,0) {$1.167P$};
    \node[above,sloped] at (2,1.5) {$-0.208P$}; \node[above,sloped] at (6,1.5) {$-1.458P$};
    \draw[load] (0,-1.4) -- (0,-0.6) node[pos=0,below] {$0.125P$}; \draw[load] (8,-1.4) -- (8,-0.7) node[pos=0,below] {$0.875P$};
    \draw[dim] (-1.6,0) -- (-1.6,3) node[midway,left] {$0.75l$};
    \node[below] at (4,-2.0) {(b) (real) internal forces $N_i$};
  \end{scope}
  % (c)
  \begin{scope}[yshift=-7.5cm]
    \trussgeo \pic at (A) {pin}; \pic at (C) {roller};
    \draw[load] (8.1,0) -- (9.1,0) node[above] {1};
    \node[above] at (2,3) {$0$}; \node[left] at (0,1.5) {$0$}; \node[right] at (4,1.5) {$0$}; \node[below] at (2,0) {$-1$}; \node[below] at (6,0) {$-1$};
    \node[above,sloped] at (2,1.5) {$0$}; \node[above,sloped] at (6,1.5) {$0$};
    \node[below] at (4,-1.5) {(c) virtual forces $n_i$ (due to $\delta P_{D1} = 1$)};
  \end{scope}
  % (d)
  \begin{scope}[xshift=12cm,yshift=-7.5cm]
    \trussgeo \pic at (A) {pin}; \pic at (C) {roller};
    \draw[load] (E) -- ++(-36.87:1.0) node[below right] {1}; \draw[load] (B) -- ++(143.13:1.0) node[above left] {1};
    \node[above] at (2,3) {$-0.8$}; \node[left] at (0,1.5) {$-0.6$}; \node[right] at (4,1.5) {$-0.6$}; \node[below] at (2,0) {$-0.8$}; \node[below] at (6,0) {$0$};
    \node[above,sloped] at (2.9,2.2) {$1$}; \node[above,sloped] at (6,1.5) {$0$};
    \node[below] at (4,-1.5) {(d) virtual forces $n_i$ (due to $\delta P_{D2} = 1$)};
  \end{scope}
\end{tikzpicture}

\end{fig}
\begin{gather*}
  \dU^* = \sum_{i=1}^{elem}n_i\delta_i\quad\text{where } \delta_i = \frac{N_i}{(EA/L)_i} = \frac{N_iL_i}{E_iA_i}\ \text{(real axial deformation)}\\
  \dW^* = \delta P_i\Delta_i\ \text{(no summation)}\\
  \text{PVF: } \dW^* = \dU^*
\end{gather*}
Calculation of $D_1$:
\begin{align*}
  (1)D_1 &= (-1)\frac{1.167P(L)}{EA} + (-1)\frac{1.167P(L)}{EA} + (0)(0) + (0)(0) + (0)\frac{-P(L)}{EA}\\
         &\quad + (0)\frac{-1.458P(1.25L)}{E(1.25A)} + (0)\frac{-0.208P(1.25L)}{E(1.25A)} = -2.334\frac{PL}{EA}
  \;\rightarrow\; D_1 = -2.334\frac{PL}{EA}
\end{align*}
Calculation of $D_2$:
\begin{align*}
  (1)D_2 &= (-0.8)\frac{1.167P(L)}{EA} + (0)\frac{1.167P(L)}{EA} + (-0.6)(0) + (-0.6)(0) + (-0.8)\frac{-P(L)}{EA}\\
         &\quad + (0)\frac{-1.458P(1.25L)}{E(1.25A)} + (1)\frac{-0.208P(1.25L)}{E(1.25A)} = -0.342\frac{PL}{EA}
  \;\rightarrow\; D_2 = -0.342\frac{PL}{EA}
\end{align*}
\end{example}

\begin{example}[Deflection of truss due to temperature/fabrication errors]
Consider the above truss but \emph{without applied loads}. Assume member 2 has a temperature rise of 30 degrees and members 5 and 6 have a temperature rise of 20 degrees. Also, assume the coefficient of thermal expansion $\alpha = 0.5\times10^{-6}/{}^\circ C$. Calculate the displacements $D_1$ and $D_2$.
\begin{gather*}
  \text{PVF: } \dW^* = \dU^*\\
  \dW^* = \delta P_i\Delta_i\quad\text{and}\quad \dU^* = \sum_{i=1}^{elem}n_i\delta_i
\end{gather*}
where $\delta = \Delta L = (\alpha\Delta T)L$ for temperature change and $\delta = \Delta L =$ fabrication error. Therefore, in this example
\begin{align*}
  \delta_2 &= \alpha\Delta T_2L_2 = (0.5\times10^{-5})(30)(L) = 15\times10^{-5}L\\
  \delta_5 &= \alpha\Delta T_5L_5 = (0.5\times10^{-5})(20)(L) = 10\times10^{-5}L\\
  \delta_6 &= \alpha\Delta T_6L_6 = (0.5\times10^{-5})(20)(1.25L) = 12.5\times10^{-5}L
\end{align*}
Calculation of $D_1$:
\begin{gather*}
  (1)(D_1) = n_2\delta_2 + n_5\delta_5 + n_6\delta_6 = (-1)(15\times10^{-5}L) + (0)(10\times10^{-5}L) + (0)(12.5\times10^{-5}L)\\
  \rightarrow\; D_1 = -15\times10^{-5}L
\end{gather*}
Calculation of $D_2$:
\begin{gather*}
  (1)(D_2) = n_2\delta_2 + n_5\delta_5 + n_6\delta_6 = (0)(15\times10^{-5}L) + (-0.8)(10\times10^{-5}L) + (0)(12.5\times10^{-5}L)\\
  \rightarrow\; D_2 = -8\times10^{-5}L
\end{gather*}
(The virtual forces $n_i$ due to $\delta P_{D1} = 1$ and $\delta P_{D2} = 1$ are those shown in (c) and (d) of the previous figure.)
\end{example}

\begin{example}[Deflection of frame due to support settlement]
\begin{equation*}
  \dW^* = (\delta P)(\Delta) + \sum_{Support}(\delta R_x\Delta_x + \delta R_y\Delta_y + \delta M\theta)
\end{equation*}
Consider the following frame with $E = 2\times10^8$~kN/m$^2$ and $I = 6\times10^{-4}$~m$^4$. Assume node $E$ has a horizontal settlement of 0.02~m to the right and node $A$ has a downward settlement of 0.03~m. Calculate the horizontal displacement of node $B$.
\begin{fig}
  % Portal frame with support settlements and a horizontal load at B
\begin{tikzpicture}[scale=0.7,font=\small]
  \draw[member] (0,0) -- (0,5) -- (10,5) -- (10,0);
  \pic at (0,0) {pin}; \pic at (10,0) {pin}; \node[node dot] at (0,0) {}; \node[node dot] at (10,0) {}; \node[node dot] at (5,5) {};
  \node[left] at (-0.3,0.1) {$A$}; \node[above left] at (0,5.05) {$B$}; \node[above] at (5,5.1) {$C$}; \node[above right] at (10,5.05) {$D$}; \node[right] at (10.3,0.1) {$E$};
  \draw[load] (-2.5,5) -- (-0.15,5) node[pos=0,left] {100 kN};
  \node[below] at (2.5,4.9) {5 m}; \node[below] at (7.5,4.9) {5 m}; \node[left] at (-0.1,2.5) {5 m};
\end{tikzpicture}

\end{fig}
\begin{fig}
  % Real moments M (100 kN) and virtual moments m (unit load) for the portal frame
\begin{tikzpicture}[scale=0.48,font=\scriptsize]
  \newcommand{\framemom}[5]{% #1 load label, #2 corner moment, #3 mid moment label, #4 reaction H, #5 reaction V
    \draw[thinmember] (0,0) -- (0,5) -- (10,5) -- (10,0);
    % column AB: triangular from 0 at A to #2 at B (drawn to the left)
    \draw (0,0) -- (-1.5,5) -- (0,5);
    \draw (0,5) -- (0,6.5) -- (10,3.5) -- (10,5);
    \draw (10,0) -- (11.5,5) -- (10,5);
    \node[above] at (-0.2,6.5) {#2}; \node[left] at (-1.5,4.7) {#2}; \node[right] at (10.4,3.5) {#2}; \node[right] at (11.5,4.7) {#2};
    \draw[load] (-3.0,5) -- (-1.6,5) node[pos=0,left] {#1};
    \draw[load] (-0.3,-0.9) -- (-1.5,-0.9) node[pos=0.5,below] {#4}; \draw[load] (0,-1.3) -- (0,-2.3) node[below] {#4};
    \draw[load] (11.5,-0.9) -- (10.3,-0.9) node[pos=0.5,below] {#4}; \draw[load] (10,-2.3) -- (10,-1.3) node[pos=0,below] {#4};
    \node at (5,2.5) {\large #3};}
  \begin{scope} \framemom{100}{250}{$M$}{50}{50} \node[below] at (5,-3.2) {(a) real moments}; \end{scope}
  \begin{scope}[xshift=17cm] \framemom{1}{2.5}{$m$}{0.5}{0.5} \node[below] at (5,-3.2) {(b) virtual moments due to $\delta P = 1$}; \end{scope}
\end{tikzpicture}

\end{fig}
\begin{gather*}
  \dW^* = 1\cdot u_B + (-0.5)(-0.03) + (-0.5)(0.02) = u_B + 0.005\\
  \dU^* = \int\frac{Mm\,dx}{EI} = \frac{1041.67}{EI} + \frac{2083.33}{EI} + \frac{1041.67}{EI} = \frac{4166.7}{EI} = \frac{4166.7}{12\times10^4} = 0.0347\\
  \text{(can be calculated using Table~\ref{tab:integral})}\\
  \dW^* = \dU^* \;\Rightarrow\; u_B = 0.0297\text{ m}
\end{gather*}
\end{example}

\clearpage
\begin{table}[H]
  \centering
  \caption{Values of the integral $\int M_uM\,dl$ (Ghali and Neville 1997, \emph{Structural Analysis}).}
  \label{tab:integral}
  % Values of the integral \int M_u M dl (after Ghali & Neville 1997)
% Small shape macros (M diagrams over length l)
\newcommand{\shpRect}[1]{\begin{tikzpicture}[scale=0.55,baseline=(current bounding box.center),font=\tiny]
  \draw[thick] (0,0) rectangle (2.4,0.7); \node[right] at (2.4,0.35) {$#1$}; \draw[dim] (0,-0.35) -- (2.4,-0.35) node[midway,fill=white,inner sep=0.5pt] {$l$};\end{tikzpicture}}
\newcommand{\shpTriR}[1]{\begin{tikzpicture}[scale=0.55,baseline=(current bounding box.center),font=\tiny]
  \draw[thick] (0,0) -- (2.4,0) -- (2.4,0.8) -- cycle; \node[right] at (2.4,0.4) {$#1$}; \draw[dim] (0,-0.35) -- (2.4,-0.35) node[midway,fill=white,inner sep=0.5pt] {$l$};\end{tikzpicture}}
\newcommand{\shpTriL}[1]{\begin{tikzpicture}[scale=0.55,baseline=(current bounding box.center),font=\tiny]
  \draw[thick] (0,0) -- (2.4,0) -- (0,0.8) -- cycle; \node[left] at (0,0.4) {$#1$}; \draw[dim] (0,-0.35) -- (2.4,-0.35) node[midway,fill=white,inner sep=0.5pt] {$l$};\end{tikzpicture}}
\newcommand{\shpTrap}[2]{\begin{tikzpicture}[scale=0.55,baseline=(current bounding box.center),font=\tiny]
  \draw[thick] (0,0) -- (2.4,0) -- (2.4,1.0) -- (0,0.45) -- cycle; \node[left] at (0,0.25) {$#1$}; \node[right] at (2.4,0.6) {$#2$}; \draw[dim] (0,-0.35) -- (2.4,-0.35) node[midway,fill=white,inner sep=0.5pt] {$l$};\end{tikzpicture}}
\newcommand{\shpPeak}[1]{\begin{tikzpicture}[scale=0.55,baseline=(current bounding box.center),font=\tiny]
  \draw[thick] (0,0) -- (2.4,0) -- (0.9,0.9) -- cycle; \draw (0.9,0) -- (0.9,0.9); \node[right] at (0.9,0.45) {$#1$};
  \draw[dim] (0,1.2) -- (0.9,1.2) node[midway,fill=white,inner sep=0.5pt] {$\alpha l$}; \draw[dim] (0.9,1.2) -- (2.4,1.2) node[midway,fill=white,inner sep=0.5pt] {$\beta l$};
  \draw[dim] (0,-0.35) -- (2.4,-0.35) node[midway,fill=white,inner sep=0.5pt] {$l$};\end{tikzpicture}}
\newcommand{\shpPara}[1]{\begin{tikzpicture}[scale=0.55,baseline=(current bounding box.center),font=\tiny]
  \draw[thick] (0,0) -- (2.4,0); \draw[thick,domain=0:2.4,samples=30] plot (\x,{0.9*(1-((\x-1.2)/1.2)^2)}); \draw (1.2,0) -- (1.2,0.9); \node[right] at (1.2,0.45) {$#1$}; \node at (2.6,0.9) {$*$};
  \draw[dim] (0,-0.35) -- (2.4,-0.35) node[midway,fill=white,inner sep=0.5pt] {$l$};\end{tikzpicture}}
\newcommand{\shpParaA}[1]{\begin{tikzpicture}[scale=0.55,baseline=(current bounding box.center),font=\tiny]
  \draw[thick] (0,0) -- (2.4,0) -- (2.4,0.9); \draw[thick,domain=0:2.4,samples=30] plot (\x,{0.9*(\x/2.4)^2}); \draw[thin] (-0.6,0) -- (0.8,0); \node[right] at (2.4,0.45) {$#1$}; \node at (2.6,1.0) {$*$};
  \node[above,font=\tiny] at (0.1,0.05) {tangent}; \draw[dim] (0,-0.35) -- (2.4,-0.35) node[midway,fill=white,inner sep=0.5pt] {$l$};\end{tikzpicture}}
\newcommand{\shpParaB}[1]{\begin{tikzpicture}[scale=0.55,baseline=(current bounding box.center),font=\tiny]
  \draw[thick] (0,0) -- (2.4,0) -- (2.4,0.9); \draw[thick,domain=0:2.4,samples=30] plot (\x,{0.9*(1-(1-\x/2.4)^2)}); \draw[thin] (1.4,0.9) -- (3.0,0.9); \node[right] at (2.4,0.45) {$#1$}; \node at (2.6,1.15) {$*$};
  \node[above,font=\tiny] at (0.5,0.55) {tangent}; \draw[dim] (0,-0.35) -- (2.4,-0.35) node[midway,fill=white,inner sep=0.5pt] {$l$};\end{tikzpicture}}
{\footnotesize
\renewcommand{\arraystretch}{1.6}
\begin{tabular}{|c|c|c|c|c|c|}\hline
  \diagbox[width=2.6cm]{$M$}{$M_u$} & \shpRect{b} & \shpTriR{b} & \shpTriL{b} & \shpTrap{b_1}{b_2} & \shpPeak{b}\\\hline
  \shpRect{a} & $abl$ & $\frac12abl$ & $\frac12abl$ & $\frac{al}{2}(b_1+b_2)$ & $\frac12abl$\\\hline
  \shpTrap{a_1}{a_2} & $\frac{bl}{2}(a_1+a_2)$ & $\frac{bl}{6}(a_1+2a_2)$ & $\frac{bl}{6}(2a_1+a_2)$ & \makecell{$\frac l6(2a_1b_1+a_1b_2+$\\$a_2b_1+2a_2b_2)$} & \makecell{$\frac{bl}{6}[(1+\beta)a_1$\\$+(1+\alpha)a_2]$}\\\hline
  \shpPeak{a} & $\frac12abl$ & $\frac{abl}{6}(1+\alpha)$ & $\frac{abl}{6}(1+\beta)$ & \makecell{$\frac{al}{6}[(1+\beta)b_1+$\\$(1+\alpha)b_2]$} & $\frac13abl$\\\hline
  \shpPara{a} & $\frac23abl$ & $\frac13abl$ & $\frac13abl$ & $\frac13al(b_1+b_2)$ & $\frac{abl}{3}(1+\alpha\beta)$\\\hline
  \shpParaA{a} & $\frac13abl$ & $\frac14abl$ & $\frac1{12}abl$ & $\frac{al}{12}(b_1+3b_2)$ & $\frac{abl}{12}(1+\alpha+\alpha^2)$\\\hline
  \shpParaB{a} & $\frac23abl$ & $\frac5{12}abl$ & $\frac14abl$ & $\frac{al}{12}(3b_1+5b_2)$ & $\frac{abl}{12}(5-\beta-\beta^2)$\\\hline
\end{tabular}\\[2pt]
{\scriptsize $*$ Second-degree parabola.}}

\end{table}
